% ============================================================================
% บทที่ 4 — จัดตารางและอ่านผลลัพธ์
% ============================================================================
\mychapter{จัดตารางและอ่านผลลัพธ์}

\noindent บทนี้อธิบายการจัดตารางสอบ การอ่านผลลัพธ์จากหน้าต่าง ๆ และการค้นหาวิชา

\section{การจัดตารางสอบ}

\begin{steps}
  \item ไปที่หน้า \textbf{ตรวจสอบข้อมูล} และตรวจสอบรายการข้อควรตรวจสอบก่อน
        (ถ้ามีปัญหา ควรแก้ก่อน)
  \item กดปุ่ม \texttt{จัดตาราง} ที่มุมขวาล่าง
\end{steps}

\begin{expected}
ระบบเริ่มคำนวณตารางสอบ ปุ่ม \texttt{จัดตาราง} เปลี่ยนเป็นปุ่ม \texttt{ยกเลิก}
เมื่อจัดเสร็จ ระบบจะแสดงหน้าภาพรวม พร้อมจำนวนวิชาที่จัดแล้วและยังไม่ได้จัด
สำหรับข้อมูลตัวอย่าง กลางภาคและปลายภาคจัดครบ \textbf{11/11} วิชา
\end{expected}

จำนวน \texttt{11/11} นับเฉพาะรายวิชาในไฟล์ \texttt{จัดอัตโนมัติ}
ที่ต้องให้ระบบจัดสอบในประเภทการสอบนั้น ข้อมูลตัวอย่างมี 12 รายวิชาแบบจัดอัตโนมัติ
แต่มีหนึ่งวิชาที่มีเงื่อนไข \texttt{ไม่มีสอบกลางภาค} และหนึ่งวิชาที่มีเงื่อนไข
\texttt{ไม่มีสอบปลายภาค} (ดูไฟล์ \file{rules.csv} ในบทที่ 2)
จึงเหลือ 11 วิชาที่ต้องจัดในแต่ละประเภท ส่วนรายวิชาในไฟล์ \texttt{จัดไปแล้ว}
ไม่ถูกนับในจำนวนนี้เลย

\begin{caution}
การกด \texttt{ยกเลิก} จะหยุดการคำนวณและคงผลลัพธ์เดิมที่เคยจัดไว้
\end{caution}

\manualfigure{docs/usage-report/screenshots/figures/F07.png}{0.95\textwidth}{%
  หน้าภาพรวมหลังจัดตาราง}

\section{การอ่านหน้าภาพรวม}

หน้าภาพรวมแสดงข้อมูลสรุปของตารางสอบ:

\begin{description}
  \item[\textbf{กลางภาค} / \textbf{ปลายภาค}] จำนวนวิชาที่จัดแล้ว จากที่ต้องจัด
        และจำนวนที่ยังไม่ได้จัด
  \item[\textbf{ต้นทางกำหนด}] จำนวนวิชาที่มีเวลาจากไฟล์ต้นทาง
  \item[\textbf{ล็อกเอง}] จำนวนวิชาที่ผู้ใช้ล็อกเวลาเอง
\end{description}

\section{การสลับประเภทการสอบ}

หน้ารายวิชา ปฏิทิน และกลุ่มนักศึกษา จะแสดงข้อมูลของประเภทการสอบหนึ่งประเภท
ในแต่ละครั้ง ให้เลือก \texttt{กลางภาค} หรือ \texttt{ปลายภาค} ที่มุมขวาบนของหน้า
ก่อนอ่านจำนวนวิชา

\needspace{9\baselineskip}
\section{การค้นหา กรอง และเปิดรายละเอียดวิชา}

ในหน้า \textbf{รายวิชา}:

\begin{steps}
  \item ใช้ช่องค้นหา \textbf{รหัสวิชา ชื่อ หรือกลุ่ม} เพื่อค้นหาวิชา
  \item ใช้ช่องกรองประเภทเพื่อเลือก \texttt{ทุกวิชา}, \texttt{จัดอัตโนมัติ} หรือ \texttt{จัดไปแล้ว}
  \item ใช้ช่องกรองสถานะเพื่อเลือก \texttt{ทุกสถานะ}, \texttt{จัดแล้ว},
        \texttt{ยังไม่ได้จัดตาราง} หรือ \texttt{ไม่จัดส่วนกลาง}
  \item เรียงลำดับตาม \texttt{รหัสวิชา} หรือ \texttt{วันสอบ}
  \item กดปุ่ม \texttt{ดูรายละเอียด} ที่อยู่ขวาของแถววิชาเพื่อเปิดหน้าต่างรายละเอียด
        ซึ่งแสดงรหัสวิชา ชื่อวิชา ประเภทการสอบ กลุ่มนักศึกษา
        แหล่งที่มาของเวลา และเวลาสอบปัจจุบัน
\end{steps}

\manualfigure{docs/usage-report/screenshots/crops/F08.png}{0.95\textwidth}{%
  หน้ารายวิชา แสดงช่องค้นหา ตัวกรอง และปุ่มดูรายละเอียดของแต่ละแถว}

\manualfigure{docs/usage-report/screenshots/crops/F09.png}{0.95\textwidth}{%
  ผลการค้นหาวิชา \course{061000001}}

\begin{note}
หากรายการวิชามีจำนวนมาก ระบบจะแสดงผลบางส่วนก่อนและมีปุ่ม \texttt{โหลดเพิ่ม}
ให้กดเพื่อแสดงรายการถัดไป
\end{note}

\section{การอ่านปฏิทิน}

หน้า \textbf{ปฏิทิน} แสดงตารางสอบรายสัปดาห์ วันเสาร์–อาทิตย์และวันหยุดมีป้ายกำกับ
วิชาที่สอบเต็มวันจะแสดงแถบเต็มความสูง และกดที่ปุ่มวิชาเพื่อเปิดรายละเอียดได้

\begin{expected}
หากวันหนึ่งมีวิชามากกว่า 5 วิชา จะแสดงปุ่ม \texttt{+N เพิ่มเติม} ให้กดเพื่อขยาย
\end{expected}

\manualfigure{docs/usage-report/screenshots/figures/F10.png}{0.95\textwidth}{%
  หน้าปฏิทินแสดงตารางสอบรายสัปดาห์}

\section{การอ่านตารางกลุ่มนักศึกษา}

หน้า \textbf{กลุ่มนักศึกษา} แสดงตารางสอบของแต่ละกลุ่ม ใช้ช่องค้นหา
\textbf{รหัสกลุ่ม} เช่น \texttt{DEMO-A} แล้วเลือกกลุ่มที่ต้องการ

\begin{expected}
ตารางแสดงรายวิชา วันสอบ และสถานะของกลุ่มที่เลือก
หากกลุ่มมีหลายวิชาในวันเดียวกัน จะแสดงข้อความ \textbf{วันที่สอบมากกว่าหนึ่งวิชา}
และหากมีวิชาที่ยังไม่ได้จัด จะแสดง \texttt{ยังไม่จัด}
\end{expected}

\manualfigure{docs/usage-report/screenshots/figures/F11.png}{0.95\textwidth}{%
  ตารางกลุ่มนักศึกษา DEMO-A}

\begin{note}
การมีหลายวิชาในวันเดียวกันไม่ได้หมายความว่าวิชาเหล่านั้นสอบพร้อมกันเสมอไป
ต้องตรวจสอบช่วงเวลาให้ละเอียด
\end{note}

\seealso{บทที่ 5 สำหรับการแก้ไขเวลาสอบ และบทที่ 6 สำหรับการตรวจสอบปัญหา}
